\documentclass[11pt]{article}

\usepackage[left=0.6 in,top=0.3in,right=0.6in,bottom=0.3in]{geometry}
\usepackage[parfill]{parskip} % Remove paragraph indentation

\usepackage{enumitem} \usepackage{hyperref} \usepackage{amsmath, amsfonts}

\begin{document} \thispagestyle{empty}

\begin{itemize} \setlength\itemsep{-.4em}
	\item Utilized YOLO v5 to develop an
	      object-detection model in PyTorch that achieved a 80\% mAP evaluation on an
	      in-house thermal-image vehicle dataset for a R\&D project.

	\item Designed an anomaly-detection model in PyTorch using an LSTM
	      variational-autoencoder to detect component failure in a pulse-power x-ray
	      machine. The model accurately predicted the location of a component failure
	      with 87\% accuracy and cut machine repair times by 50\%.

	\item Applied the U-Net architecture with PyTorch to develop a deblurring
	      model for a pulse-powered x-ray machine. The model was applied by the
	      machine operators and reduced machine configuration time by 30\%.

	\item Enhanced the scalability of models initially prototyped by our team of
	      Analysts and Data Scientists for training on our internal high-performance
	      computing (HPC) cluster. My role involved refactoring the existing code to
	      enable efficient execution in a distributed and parallel setting.

	\item Mentored a summer intern in the development and execution of models on
	      a HPC cluster, involving multi-node, multi-GPU setups and Slurm resource
	      management.

	\item Developed and deployed a Django web interface with Docker Compose,
	      enabling non-technical staff to add nodes and connections to the company's
	      Neo4j graph database, without knowledge of the query syntax.

	\item Developed interactive dashboards using Streamlit to give
	      non-technical staff the abilty to interact with the capabilities designed
	      by our team of Data Scientists.
\end{itemize}

\end{document}

